\begin{frame}[shrink]
	\frametitle{Retropropagação através do tempo (BPTT)}
	\small
	\begin{columns}[T]
		\column{.55\linewidth}
		\resizebox{!}{3.3cm}{				\begin{tikzpicture}[
	box/.style={draw, rectangle, rounded corners=3pt,
		minimum width=1.2cm, minimum height=1.0cm,
		fill=gray!15, font=\small},
	fwd/.style={->, thick, blue!65!black},
	bwd/.style={->, thick, dashed, red!65!black},
	lbl/.style={font=\footnotesize}
	]
	% Forma compacta
	\node[box] (rc) {$\mathbf{W}$};
	\coordinate (xc) at ([yshift=0.9cm]rc.north);
	\node[lbl, above] at (xc) {$x_t$};
	\draw[fwd] (xc) -- (rc.north);
	\coordinate (sc) at ([yshift=-0.9cm]rc.south);
	\node[lbl, below] at (sc) {$s_t$};
	\draw[fwd] (rc.south) -- (sc);
	\draw[fwd] (rc.north east)
	to[out=50,in=-50,looseness=3.5]
	node[lbl, right] {$h_t$}
	(rc.south east);
	\node[lbl, below=2.4cm of rc, align=center] {Forma\\compacta};
	
	% Seta de equivalência
	\node[right=1.6cm of rc, font=\Large] {$\Rightarrow$};
	
	% Forma desenrolada
	\node[box, right=3.0cm of rc] (r1) {$\mathbf{W}$};
	\node[box, right=2.2cm of r1]  (r2) {$\mathbf{W}$};
	\node[lbl, right=1.4cm of r2]  (rd) {$\cdots$};
	\node[box, right=1.4cm of rd]  (rT) {$\mathbf{W}$};
	
	% Entradas
	\foreach \nd/\xi in {r1/1, r2/2, rT/T}{
		\coordinate (inp\nd) at ([yshift=0.9cm]\nd.north);
		\node[lbl, above] at (inp\nd) {$x_{\xi}$};
		\draw[fwd] (inp\nd) -- (\nd.north);
	}
	
	% Estado recorrente (forward)
	\draw[fwd] (r1.east) -- node[lbl, above] {$h_1$} (r2.west);
	\draw[fwd] (r2.east) -- node[lbl, above] {$h_2$} (rd.west);
	\draw[fwd] (rd.east) -- node[lbl, above] {$h_{T-1}$} (rT.west);
	
	% Saídas e perdas
	\foreach \nd/\xi in {r1/1, r2/2, rT/T}{
		\node[lbl, below=0.8cm of \nd] (L\xi) {$\mathcal{L}_{\xi}$};
		\draw[fwd] (\nd.south) -- (L\xi);
	}
	
	% Retropropagação do gradiente
	\draw[bwd] ([yshift=-0.3cm]r2.west) -- ([yshift=-0.3cm]r1.east);
	\draw[bwd] ([yshift=-0.3cm]rd.west) -- ([yshift=-0.3cm]r2.east);
	\draw[bwd] ([yshift=-0.3cm]rT.west) -- ([yshift=-0.3cm]rd.east);
	
	\node[lbl, below=1.6cm of r2, text=red!65!black]
	{$\longleftarrow$ gradiente retropropagado};
	\node[lbl, below=2.3cm of r2]
	{pesos $W$ compartilhados em todas as cópias};
	
\end{tikzpicture}}
		\column{.4\linewidth}
		\begin{itemize}
			\item A rede é \textbf{desenrolada} em $T$ cópias no tempo, compartilhando os mesmos pesos $W$.
			\item \textit{Forward}: o estado $h_t$ avança da esquerda para a direita.
			\item \textit{Backward}: o gradiente percorre o caminho inverso, acumulando a contribuição de cada passo:
		\end{itemize}
		\begin{equation*}
			\frac{\partial \mathcal{L}}{\partial W} = \sum_{t=1}^{T} \frac{\partial \mathcal{L}_t}{\partial W}
		\end{equation*}
	\end{columns}
	\vspace{2pt}
	\par\sloppy É neste \textit{backward}, a cada passo $t$, que o gradiente substituto entra: ele troca a derivada (nula) da função de disparo real por uma aproximação suave, deixando o gradiente fluir de volta pelos $T$ passos \cite{wu2018spatio, eshraghian2023training}.
	% [SOFTWARE] SNN -> Surrogate gradient | mesma demo de snnTreinamento.tex
	\abrirNoSoftware{SNN \textrightarrow{} Surrogate gradient}{snn.surrogate}
\end{frame}
